\documentclass[conference]{IEEEtran}
\IEEEoverridecommandlockouts
% The preceding line is only needed to identify funding in the first footnote. If that is unneeded, please comment it out.
\usepackage{cite}
\usepackage{amsmath,amssymb,amsfonts}
\usepackage{algorithmic}
\usepackage{graphicx}
\usepackage{textcomp}
\usepackage{xcolor}
\usepackage{booktabs}
\def\BibTeX{{\rm B\kern-.05em{\sc i\kern-.025em b}\kern-.08em
    T\kern-.1667em\lower.7ex\hbox{E}\kern-.125emX}}
% Let columns end naturally instead of stretching vertical space to fill the page.
\raggedbottom
\begin{document}

\title{Autonomous Drone Racing Report}

\author{\IEEEauthorblockN{1\textsuperscript{st} Matuszewski Timo}
\IEEEauthorblockA{\textit{School of Computation, Information and Technology} \\
\textit{Technical University of Munich}\\
Munich, Germany}
\and
\IEEEauthorblockN{2\textsuperscript{nd} Kenny Jeffrey}
\IEEEauthorblockA{\textit{School of Computation, Information and Technology} \\
\textit{Technical University of Munich}\\
Munich, Germany}
}

\maketitle

\begin{abstract}
We present an autonomous drone-racing controller based on Model Predictive
Contouring Control (MPCC). A quadrotor must fly through a sequence of gates in
order while avoiding obstacles, with gate and obstacle poses only revealed once
the drone is within sensor range, forcing online replanning. We embed a
geometric racing line, produced by a B-spline planner, directly into the
prediction model as a function of a progress state, so that the reference cannot
"run away" from the drone, a failure mode we observed with a time-indexed
reference-tracking Model Predictive Control (MPC). On top of the nominal MPCC we add three lightweight,
context-dependent mechanisms: a curvature-based speed cap, a near-gate
contouring boost, and a caution mode that slows the drone near unmeasured
objects. We also implemented a reinforcement-learning (RL) layer that adapts the
MPCC cost weights online, but found end-to-end training impractical because the
embedded solver is sequential; the hand-designed context heuristics distil the
same idea at negligible cost. On a fixed racing track our controller finishes
$88\%$ of runs at a $7.37\,\mathrm{s}$ average lap time, and generalises to
fully randomised tracks at $87\%$. Deployed unchanged on the physical
Crazyflie, it completed all eight test flights at an $8.70\,\mathrm{s}$ mean lap
time, confirming the sim-to-real transfer. An ablation study quantifies the
contribution of each component.
\end{abstract}

\section{Introduction}
Autonomous drone racing has become a popular benchmark for agile, high-speed
robot control, because it couples fast nonlinear dynamics, tight actuation
limits, and hard geometric constraints (the gates) in a single task
\cite{hanover2024survey}. Recent work has pushed quadrotors to champion-level
performance using both learning-based \cite{kaufmann2023champion} and
optimisation-based \cite{foehn2021timeoptimal,romero2022mpcc} approaches, and the
gap between the two is an active research direction. Beyond racing, the
underlying problem (flying a prescribed path as fast as possible while
respecting dynamics and safety limits) is representative of a broad class of
agile aerial applications such as inspection and delivery.

In this work we consider the drone-racing task posed by the course simulator: a
Crazyflie~2.1 quadrotor must pass four gates in a fixed order and reach the
finish, while avoiding four pole obstacles. The challenge is threefold. First,
the flight is \emph{time-critical}: the objective is to minimise lap time, which
drives the drone to the edge of its actuation envelope, where tracking errors
translate directly into gate-frame collisions. Second, the environment is only
\emph{partially observed}: gates and obstacles are known only at nominal poses
until the drone comes within a short sensor range, at which point the true pose
is revealed and the planned path must be corrected online. Third, the controller
must \emph{generalise}: it is ultimately graded on a realistic, hand-set track on
real hardware, so overfitting to one geometry is penalised by the sim-to-real
gap and by unseen layouts.

We address this with Model Predictive Contouring Control (MPCC)
\cite{lam2010mpcc,liniger2015racing,romero2022mpcc,krinner2024mpccpp}, which optimises progress
along a geometric reference path rather than tracking a time-parameterised
trajectory. We contribute (i) a complete MPCC realisation for the Crazyflie with
soft obstacle and gate-frame avoidance embedded in the cost, (ii) three
context-dependent adaptation mechanisms motivated by an RL weight-scheduling idea
that we found too expensive to train end-to-end, and (iii) a systematic ablation
study on both a fixed and a fully randomised track that isolates the effect of
each component.

\section{Related Work}
\emph{Optimisation-based racing.} Contouring control was introduced for machining
\cite{lam2010mpcc} and adapted to autonomous racing of scale cars
\cite{liniger2015racing}, where progress along a reference is maximised online.
Romero et al.\ transferred MPCC to quadrotors and showed near-time-optimal flight
\cite{romero2022mpcc}, and Krinner et al.\ extended it with hard safety constraints in MPCC++ \cite{krinner2024mpccpp}, while time-optimal planning approaches compute
aggressive reference trajectories offline \cite{foehn2021timeoptimal,penicka2022}.
Our work builds directly on this quadrotor MPCC line
\cite{romero2022mpcc,krinner2024mpccpp}, but keeps collision avoidance soft in the
cost rather than as MPCC++'s hard tunnel, and targets a partially observed, replanning setting rather than a known
track.

\emph{Learning-based racing and control.} End-to-end reinforcement learning has
reached champion-level performance from onboard sensing
\cite{kaufmann2023champion,song2021rl}, at the cost of large-scale training.
A complementary line keeps a model-based controller and learns only part of it:
learning the plant model \cite{kabzan2019lmpc} or, as here, the cost-function
weights \cite{zarrouki2021wmpc,zanon2021rlmpc,mehndiratta2018}. We adopt the
latter view but, finding end-to-end weight learning too costly with an embedded
solver, distil it into interpretable context rules.

\section{Problem Formulation}
We model the quadrotor as a continuous-time nonlinear system with state
$\mathbf{x}\in\mathbb{R}^{12}$ and input $\mathbf{u}\in\mathbb{R}^{4}$,
\begin{equation}
  \mathbf{x} = \big[\mathbf{p}^\top,\,\boldsymbol{\Phi}^\top,\,
  \mathbf{v}^\top,\,\boldsymbol{\omega}^\top\big]^\top,\qquad
  \mathbf{u} = \big[\boldsymbol{\Phi}_{\mathrm{cmd}}^\top,\,T\big]^\top,
  \label{eq:state}
\end{equation}
where $\mathbf{p}$ is the world position, $\boldsymbol{\Phi}=(\phi,\theta_\mathrm{p},\psi)$
the roll--pitch--yaw attitude, $\mathbf{v}$ the linear velocity,
$\boldsymbol{\omega}$ the body rates, $\boldsymbol{\Phi}_{\mathrm{cmd}}$ the
commanded attitude and $T$ the collective thrust; the continuous dynamics
$\dot{\mathbf{x}}=f_{\mathrm{q}}(\mathbf{x},\mathbf{u})$ follow the reduced
attitude-command quadrotor model of \cite{mellinger2011minsnap}. The controller
runs at $50\,\mathrm{Hz}$ and outputs the attitude command $\mathbf{u}$; the
simulator integrates the full rigid-body dynamics at $500\,\mathrm{Hz}$.

The track is an ordered set of $N_g$ gates with centre poses
$\{\mathbf{g}_i\}_{i=1}^{N_g}$ and $N_o$ cylindrical obstacles. A gate $i$
counts as \emph{passed} when the drone crosses its plane, from the front, inside
a $0.45\times0.45\,\mathrm{m}$ window. Let $g\in\{0,\dots,N_g-1,-1\}$ be the
index of the next gate to pass, initialised at $0$ and set to $-1$ once the last
gate is cleared. A run is a \emph{success} iff $g=-1$ before the horizon of
$1500$ steps ($30\,\mathrm{s}$) elapses. The drone is \emph{disabled}
(terminating the episode) on any contact with a gate or obstacle, or on leaving
the safety box $\mathcal{X}_{\mathrm{safe}}$.

Crucially, the observation exposes each gate/obstacle at a fixed \emph{nominal}
pose until the drone is within a sensor range $r_s=0.7\,\mathrm{m}$, after which
the \emph{true} (randomly perturbed) pose is revealed. The control problem is
therefore
\begin{equation}
\begin{aligned}
  \min_{\mathbf{u}(\cdot)} \;\; & t_{\mathrm{lap}}\\
  \text{s.t.}\;\; & \dot{\mathbf{x}}=f_{\mathrm{q}}(\mathbf{x},\mathbf{u}),\quad
  \mathbf{u}\in\mathcal{U},\;\mathbf{x}\in\mathcal{X}_{\mathrm{safe}},\\
  & \text{pass gates } 1,\dots,N_g \text{ in order},\\
  & \|\mathbf{p}-\mathbf{o}_j\|\ge r_j \;\;\forall j
  \text{ (collision-free)},
\end{aligned}
\label{eq:problem}
\end{equation}
solved online under partial observability of $\{\mathbf{g}_i\}$ and
$\{\mathbf{o}_j\}$.


\section{Methodology}
\subsection{System overview}
Our controller (Fig.~\ref{fig:arch}) has three stages: a \emph{planner} that
builds a smooth, collision-aware geometric racing line through the gates; an
\emph{MPCC} that optimises progress along that line while respecting dynamics,
input, and soft avoidance constraints; and a set of \emph{context-adaptation}
rules that modulate the MPCC speed target and weights from the local track
geometry. When any measured object pose changes, the planner replans in a
background thread and the new path is swapped into the MPCC.
\begin{figure}[!ht]
  \centering
  \includegraphics[width=\columnwidth]{figures/mpc_blockdiagram.pdf}
  \caption{Control architecture: the planner builds a geometric racing line,
  the MPCC optimises progress along it, and context rules modulate the speed
  target and weights. A pose change triggers a background replan.}
  \label{fig:arch}
\end{figure}

\subsection{Path planning}
For each inter-gate segment we place three fixed waypoints per gate (before,
centre, after) along the gate normal, plus $N_f=4$ free intermediate waypoints
whose positions are optimised with L-BFGS-B against a cost that penalises
proximity to obstacles and gate frames and deviation from the straight
connection. A weighted cubic B-spline is fitted through the waypoints
\cite{dierckx1993splines} and reparameterised by arc length $\theta$, giving a
$C^2$ reference curve $\mathbf{p}_d(\theta)$ with unit tangent
$\mathbf{t}(\theta)=\mathbf{p}_d'(\theta)$. Arc-length parameterisation makes the
progress state physically meaningful and decouples the path geometry from the
speed profile. Fig.~\ref{fig:plan} shows a resulting line, with the free
waypoints bowed clear of the obstacles.
The free waypoints $\{\mathbf{w}_k\}_{k=1}^{N_f}$ of each segment are placed by
minimising
\begin{equation}
  J_{\mathrm{plan}} = \sum_k \Big( w_o\!\sum_j \beta(\mathbf{w}_k,\mathcal{C}_j)
  + w_f\,\Gamma(\mathbf{w}_k) + w_d\,\|\mathbf{w}_k-\bar{\mathbf{w}}_k\|^2\Big),
  \label{eq:plancost}
\end{equation}
where $\beta$ is the obstacle barrier of \eqref{eq:barrier}, $\Gamma$ penalises
penetration of any non-target gate frame, and the last term anchors the waypoint
to the straight-line position $\bar{\mathbf{w}}_k$ so the line stays short. The
problem is solved with L-BFGS-B, warm-started from the previous solution so a
replan converges in a few iterations.
\begin{figure}[!ht]
  \centering
  \includegraphics[width=\columnwidth]{figures/planning_topdown.pdf}
  \caption{Top-down view of a planned racing line through the gates, with the
  free intermediate waypoints bowed away from the obstacles.}
  \label{fig:plan}
\end{figure}

\subsection{MPCC formulation}
Following \cite{lam2010mpcc,liniger2015racing,romero2022mpcc,krinner2024mpccpp} we augment the
quadrotor state with a \emph{progress} state $\theta$ and its rate $v_\theta$,
driven by a virtual input $a_\theta$:
\begin{equation}
  \tilde{\mathbf{x}}=[\mathbf{x}^\top,\theta,v_\theta]^\top\in\mathbb{R}^{14},\quad
  \tilde{\mathbf{u}}=[\mathbf{u}^\top,a_\theta]^\top\in\mathbb{R}^{5},
  \label{eq:aug}
\end{equation}
with $\dot\theta=v_\theta$, $\dot v_\theta=a_\theta$. At a predicted position
$\mathbf{p}$ and progress $\theta$ we define the path deviation
$\mathbf{d}=\mathbf{p}-\mathbf{p}_d(\theta)$ and split it into a \emph{lag}
(longitudinal) and a \emph{contouring} (perpendicular) error
(Fig.~\ref{fig:contouring})
\begin{equation}
  e_l = \mathbf{t}(\theta)^\top\mathbf{d},\qquad
  \mathbf{e}_c = \mathbf{d}-e_l\,\mathbf{t}(\theta),
  \label{eq:errors}
\end{equation}
\begin{figure}[!ht]
  \centering
  \includegraphics[width=0.9\columnwidth]{figures/contouring_schematic.pdf}
  \caption{Contouring formulation: the deviation $d=p-p_d(\theta)$ is split
  into the lag error $e_l$ along the path tangent $t(\theta)$ and the
  perpendicular contouring error $e_c$.}
  \label{fig:contouring}
\end{figure}
and a progress-speed error $e_v = v_\theta - v_{\mathrm{tgt}}$ that rewards fast
travel along the path. The MPCC minimises, over the horizon of $N=25$ nodes with
step $\mathrm{d}t=0.02\,\mathrm{s}$ ($T_p=0.5\,\mathrm{s}$),
\begin{equation}
\begin{aligned}
  J = \sum_{k=0}^{N}\Big(&
  \|\mathbf{e}_{c,k}\|_{q_c}^2 + q_l e_{l,k}^2 + q_v e_{v,k}^2\\
  &+ \|\boldsymbol{\Phi}_k\|_{q_\Phi}^2 + \|\boldsymbol{\omega}_k\|_{q_\omega}^2
  + \|\tilde{\mathbf{u}}_k\|_{R}^2\\
  &+ \textstyle\sum_j w_b\,\beta(\mathbf{p}_k,\mathcal{C}_j)
  + w_g\,\gamma(\mathbf{p}_k)\Big),
\end{aligned}
\label{eq:cost}
\end{equation}
expressed as a nonlinear least-squares residual with zero reference. The first
line is the core MPCC objective: track the line ($\mathbf{e}_c$), keep the
progress state glued to the drone ($e_l$), and cruise ($e_v$). We deliberately
set $q_l\!\gg\! q_v$ (here $150$ vs.\ $5$) so the progress state cannot sprint
ahead of the physical drone. The remaining terms penalise attitude, body rate,
and control effort. Collision avoidance is \emph{soft}, in contrast to the hard tunnel constraint of MPCC++ \cite{krinner2024mpccpp}: each obstacle and gate
frame is modelled as a capsule $\mathcal{C}_j$ and enters the cost through a
one-sided barrier residual
\begin{equation}
  \beta(\mathbf{p},\mathcal{C}_j)=
  \max\!\Big(0,\;1-\tfrac{\mathrm{dist}(\mathbf{p},\mathcal{C}_j)^2}{r_j^2}\Big)^2,
  \label{eq:barrier}
\end{equation}
with a large weight $w_b$; a further one-sided term $\gamma$ enforces ground
clearance. Keeping avoidance in the cost (rather than as hard constraints) means
the QP is always feasible and the gate opening stays free to fly through. State
and input box limits ($\mathbf{v}$, $v_\theta$, attitude, thrust, $a_\theta$)
are imposed as constraints, with the velocity bounds softened by slack so a
transient overspeed never renders the problem infeasible. The full set of weights
and solver settings is summarised in Table~\ref{tab:params}.

\begin{table}[t]
\centering
\caption{Main MPCC weights and solver settings from:
\texttt{mpcc/config.py}).}
\label{tab:params}
\begin{tabular}{llll}
\toprule
Symbol & Value & Symbol & Value\\
\midrule
$q_c$ (contour) & $50$ & horizon $N$ & $25$\\
$q_l$ (lag) & $150$ & step $\mathrm{d}t$ & $0.02$\,s\\
$q_v$ (progress) & $5$ & $T_p$ & $0.5$\,s\\
$q_\Phi$ (att.) & $1$ & $w_b$ (barrier) & $5000$\\
$q_\omega$ (rate) & $5$ & $w_g$ (ground) & $400$\\
$r_\Phi/r_T/r_a$ (effort) & $1/50/0.5$ & $a_{\mathrm{lat}}^{\max}$ & $20$\,m/s$^2$\\
$\beta_g$ / $\rho$ (boost) & $4$ / $0.5$\,m & $v_{\min}$ & $0.8$\,m/s\\
$\kappa_c$ / $R_c$ (caution) & $0.5$ / $1.3$\,m & $V_{\mathrm{tgt}}$ & $2.5$\,m/s\\
\bottomrule
\end{tabular}
\end{table}


\subsection{Real-time solver}
Problem \eqref{eq:cost} is transcribed with direct multiple shooting
(explicit Runge--Kutta integration) and solved with \texttt{acados}
\cite{verschueren2022acados} using a Gauss--Newton Hessian and the Real-Time
Iteration (RTI) scheme \cite{diehl2005rti}, which replaces a fully converged SQP
solve with a single SQP iteration per control tick, split into a preparation and a
feedback phase; the resulting QP is solved by \texttt{HPIPM} \cite{frison2020hpipm}
in full-condensing mode, and the symbolic model and cost are built with CasADi
\cite{andersson2019casadi}. A receding-horizon warm start shifts the previous
solution one node forward, keeping each tick near the previous optimum, so a single
iteration suffices and the optimiser converges \emph{across} control ticks rather
than within one. On a path swap, where this no longer holds, the solver is re-seeded
and re-converged with a short iteration burst. The fixed per-tick cost is what keeps
the $50\,\mathrm{Hz}$ deadline reliable, whereas a fully converged SQP would spend a
data-dependent number of iterations and risk missing it. The generated C solver is
cached by a source-and-settings hash and reused across episodes, so only the first
run pays the code-generation cost.



\subsection{Why contouring control instead of reference tracking}
Our first controller was a reference-tracking MPC that followed a
\emph{time-indexed} trajectory $\mathbf{p}_{\mathrm{ref}}(t_k)$. Because the
reference marches forward on a clock regardless of the drone's actual progress,
whenever the drone fell behind (e.g.\ after a replan or an aggressive corner) the
MPC chased a far-ahead point at full thrust and the reference effectively "ran
away", requiring a hand-tuned reference governor to hold it back. MPCC removes
this failure mode structurally: the reference is $\mathbf{p}_d(\theta)$, a
function of a \emph{state the optimiser itself controls}, so it can never outrun
the drone \cite{romero2022mpcc}. This is the central reason we adopted MPCC; it
also aligns naturally with the model-predictive path-following formulation of
\cite{faulwasser2016mpfc}.

Intuitively, $\theta$ is a virtual particle running along the racing line whose
speed the optimiser sets, with the drone tethered to it by the contouring and lag
costs. Because the far horizon nodes are evaluated at $\mathbf{p}_d(\theta)$ inside
an upcoming corner (where a small contouring error at high speed would demand
infeasible attitude), the optimiser slows $\theta$ before the apex and
accelerates out, so MPCC already brakes implicitly through corners \cite{romero2022mpcc,krinner2024mpccpp}. This only works
if the path enters the cost as the live function $\mathbf{p}_d(\theta)$;
freezing a per-tick reference point drops $\theta$ from the objective and silently
reduces the controller to point-tracking MPC, the very failure mode above. The lag
weight $q_l\!\gg\!q_v$ is the soft anchor that keeps $\theta$ glued to the drone
through the turn. Because our progress reward tracks a target speed rather than maximising raw progress, this implicit braking alone still carries too much speed into sharp turns; the explicit curvature cap of the next subsection lowers that target into corners. 

\subsection{Context-dependent adaptation}
Two properties of the track are known to the planner but not encoded in the
static weights: how sharp the upcoming corner is, and how close the drone is to
a gate. We exploit both.

\emph{Curvature speed cap.} From the arc-length spline we compute the path
curvature $\kappa(s)=\|\mathbf{p}_d''(s)\|$, smooth it over a short window, and
set a lateral-acceleration-limited speed cap
\begin{equation}
  v_{\mathrm{cap}}(s)=\mathrm{clip}\!\Big(\sqrt{a_{\mathrm{lat}}^{\max}/\kappa(s)},
  \; v_{\min},\; V_{\mathrm{tgt}}\Big),
  \label{eq:vcap}
\end{equation}
taking the tightest cap over a $1.5\,\mathrm{m}$ look-ahead so the drone brakes
\emph{before} the apex. The MPCC speed target $v_{\mathrm{tgt}}$ in $e_v$ is set
to this cap (Fig.~\ref{fig:speed}).

\begin{figure}[!ht]
  \centering
  \includegraphics[width=\columnwidth]{figures/speed_profile.pdf}
  \caption{Speed along the planned path: the planner cruises at the target
    $V_{\mathrm{tgt}}$ (trapezoidal, with accel/decel ramps), while the controller's
    curvature cap $v_{\mathrm{cap}}(s)=\sqrt{a_{\mathrm{lat}}^{\max}/\kappa(s)}$
    lowers the effective target into corners (down to the floor
    $v_{\min}$) and leaves straights at full speed.}
  \label{fig:speed}
\end{figure}

\emph{Near-gate contouring boost.} The drone tends to drift laterally exactly at
a gate and clip the frame. For any prediction node whose progress lies within a
window $\rho$ of a gate centre $\theta_{g_i}$ we multiply the contouring weight,
\begin{equation}
  q_c \leftarrow \beta_g\, q_c \quad\text{if}\quad
  \min_i|\theta-\theta_{g_i}|\le\rho,
  \label{eq:boost}
\end{equation}
($\beta_g=4$, $\rho=0.5\,\mathrm{m}$), pulling the prediction tightly onto the
gate-centred line where precision matters, while straights keep the loose, fast
baseline weighting.

\emph{Caution mode.} While the true pose of a gate/obstacle is still unknown, the
target speed is scaled down toward a fraction $\kappa_c$ of $V_{\mathrm{tgt}}$ as
the drone closes within a radius $R_c$ of it, so that it is already slow when the
true pose is revealed at $r_s$ and the path is corrected.

\subsection{Online replanning}
A change in any measured gate/obstacle pose (detected when the true pose is
revealed within $r_s$) triggers a full-track replan in a background thread, so
the $50\,\mathrm{Hz}$ loop never stalls. When the new path is ready it is swapped
in, the embedded spline coefficients are refreshed, and the progress state is
re-localised by projecting the drone onto the new arc length. A gate-progress
barrier additionally prevents $\theta$ from advancing past the current target
gate until the pass is confirmed, so the optimiser cannot "skip" a gate.
Concretely, the per-node progress anchor is clamped to
\begin{equation}
  \theta_k \le \theta_{g} + \Delta_{\mathrm{lead}},\qquad
  \Delta_{\mathrm{lead}}=0.4\,\mathrm{m},
  \label{eq:gatebar}
\end{equation}
where $\theta_{g}$ is the arc length of the current target gate; the bound only
advances to the next gate once the environment confirms the pass.

\subsection{Learned weight adaptation (attempted)}
The three rules above are hand-designed instances of a more general idea:
schedule the MPCC cost weights online from context. We therefore implemented a
Weights-varying MPC in the spirit of \cite{zarrouki2021wmpc}, where a policy
outputs, every few ticks, bounded multipliers on grouped baseline weights
\begin{equation}
  m = \exp\!\big(\ln(4)\cdot\tanh(a)\big)\in[\tfrac14,4],\quad a\in\mathbb{R}^{6},
  \label{eq:mult}
\end{equation}
(so $a=0$ recovers the exact baseline), from a $27$-dimensional feature vector
(tracking errors, speed, look-ahead curvature, and relative gate/obstacle
geometry). We trained it with PPO \cite{schulman2017ppo} with the MPCC
\emph{inside} the environment loop, and a reward combining path progress, gate
bonuses, and crash/effort penalties; this places our work in the family of
RL-tuned MPC \cite{mehndiratta2018,zanon2021rlmpc}. In practice, training was
impractical on our hardware: because \texttt{acados} is a sequential C solver,
the environment cannot be vectorised like a pure simulator, so each rollout runs
one MPCC at a time and data collection is slow. Rather than pay that training
cost, we distilled the same context-dependent intuition into the deterministic
curvature-, gate-, and caution rules of \eqref{eq:vcap}--\eqref{eq:boost}, which
require no training and are trivially interpretable. The policy interface remains
in the code (disabled by default) for future work.


\section{Results}
\subsection{Experimental setup}
We evaluate in the course simulator (\texttt{crazyflow}, a JAX/MuJoCo-MJX
Crazyflie~2.1 model). Two settings are used: \emph{Level~2}, a single fixed,
hand-designed four-gate track with per-episode position/attitude jitter
($\pm0.15\,\mathrm{m}$ on gates/obstacles, plus inertial and actuation
disturbances); and \emph{Level~3}, identical disturbances but with the whole
track geometry randomised every episode. All numbers below are over $n=100$
episodes at a \emph{fixed} seed ($\mathrm{seed}=0$): because the seed is held
constant across configurations, every variant is evaluated on the \emph{same}
sequence of disturbances (and, on Level~3, the same sequence of random
tracks), so the comparisons are \emph{paired}. We treat Level~2 as the
\emph{development} track (on which the controller was tuned) and Level~3 as a
\emph{held-out} test of generalisation. A run is a success if all four gates are
passed. We report success rate, the mean/median lap time over \emph{successful}
runs, and a breakdown of failures by the nearest physical object at the crash.

\subsection{Nominal performance}
Table~\ref{tab:main} reports the full controller. On the fixed track it finishes
$88\%$ of runs at $7.37\,\mathrm{s}$ mean lap time, and, importantly, it
retains $87\%$ on the fully randomised track, indicating that the method does not
overfit the single Level~2 geometry. The mean lap time on Level~3
($8.51\,\mathrm{s}$) is higher than on Level~2, but the two are not directly
comparable: the randomised Level~3 tracks are on average longer than the fixed
Level~2 layout, so the difference chiefly reflects track length rather than a
slower or more conservative controller. Failures are dominated by gate-frame clips at the tight
$180^\circ$ reversal (gate~2) and at the final gate, where a late pose reveal
leaves the least room to correct.

\begin{table}[t]
\centering
\caption{Nominal performance ($n=100$, per level).}
\label{tab:main}
\begin{tabular}{lccc}
\toprule
Track & Success & Mean time [s] & Median [s]\\
\midrule
Level~2 (fixed) & $88.0\%$ & $7.37$ & $7.30$\\
Level~3 (randomised) & $87.0\%$ & $8.51$ & $8.50$\\
\bottomrule
\end{tabular}
\end{table}

\subsection{Real-world deployment}
The ultimate test is the physical Crazyflie. We deployed the \emph{exact} controller
(config \texttt{final.toml}, $V_{\mathrm{tgt}}=2.5\,\mathrm{m/s}$, no re-tuning) on
the real drone with a motion-capture state estimate. Across eight consecutive
flights the drone completed the track \emph{every time}, with lap times of
$\{8.20, 8.53, 8.54, 8.71, 8.72, 8.79, 8.84, 9.24\}\,\mathrm{s}$
(Table~\ref{tab:real}). The mean of $8.70\,\mathrm{s}$ lies within
$0.2\,\mathrm{s}$ of the Level-3 simulation mean ($8.51\,\mathrm{s}$), a small
sim-to-real gap: the physical laps are marginally slower and more variable,
consistent with unmodelled aerodynamics, battery sag, and estimator latency. The
RTI solve held the $50\,\mathrm{Hz}$ deadline, with only occasional overruns of a
few milliseconds ($\le 10\,\mathrm{ms}$) that did not affect completion. The controller thus
transfers to hardware without modification, validating the simulation results.

\begin{table}[t]
\centering
\caption{Real-world deployment on the physical Crazyflie
(\texttt{final.toml}, $V_{\mathrm{tgt}}=2.5\,\mathrm{m/s}$): eight consecutive
completed flights.}
\label{tab:real}
\begin{tabular}{lc}
\toprule
Metric & Value\\
\midrule
Completed flights & $8/8$\\
Mean lap time & $8.70\,\mathrm{s}$\\
Median lap time & $8.71\,\mathrm{s}$\\
Best / worst & $8.20$ / $9.24\,\mathrm{s}$\\
Std.\ deviation & $0.30\,\mathrm{s}$\\
Level-3 simulation reference & $8.51\,\mathrm{s}$\\
\bottomrule
\end{tabular}
\end{table}

\subsection{Ablation study}
To quantify each component we perform a leave-one-out ablation (remove one
mechanism, keep the rest) and an all-off reference (raw MPCC), all on Level~2
(Table~\ref{tab:ablation}).

\begin{table}[t]
\centering
\caption{Ablation on Level~2 ($n=100$). "Frame" = gate-frame clips,
"Coll." = total collisions.}
\label{tab:ablation}
\begin{tabular}{lcccc}
\toprule
Configuration & Success & Mean [s] & Frame & Coll.\\
\midrule
Full (all on)          & $88.0\%$ & $7.37$ & $4$  & $5$\\
$-$ Gate-track boost   & $83.0\%$ & $7.46$ & $10$ & $8$\\
$-$ Curvature cap      & $88.0\%$ & $7.15$ & $6$  & $7$\\
$-$ Caution mode       & $79.0\%$ & $6.13$ & $15$ & $17$\\
All off (raw MPCC)     & $62.0\%$ & $5.97$ & $25$ & $27$\\
\bottomrule
\end{tabular}
\end{table}
\begin{figure}[!ht]
  \centering
  \includegraphics[width=\columnwidth]{figures/ablation_bars.pdf}
  \caption{Leave-one-out ablation on Level~2 ($n=100$): success rate (bars,
  left) and mean lap time (line, right). Caution is the dominant robustness
  lever; removing all rules yields the fast but fragile raw MPCC.}
  \label{fig:ablation}
\end{figure}
The results (Fig.~\ref{fig:ablation}) isolate three distinct roles.
\emph{Caution} is the dominant robustness lever: removing it drops success by $9$ points and nearly triples
collisions ($5\!\to\!17$), because the drone enters a blind gate at full speed
and cannot absorb the replan when the true pose appears at $r_s=0.7\,\mathrm{m}$;
in exchange it saves $1.24\,\mathrm{s}$, making explicit the speed/robustness
trade-off. The \emph{gate-track boost} is a precision lever: removing it roughly
doubles gate-frame clips ($4\!\to\!10$) and costs $5$ points of success at
essentially no change in lap time, confirming that tightening $q_c$ near gates
buys frame clearance almost for free. The \emph{curvature cap} is neutral in
success at our conservative $V_{\mathrm{tgt}}$ and is in fact marginally faster
when removed; its value is as a \emph{speed-scaling safety margin} that becomes
essential once $V_{\mathrm{tgt}}$ is raised, keeping a fast configuration from
overshooting the reversal. Stacking all three converts the raw MPCC
($62\%$/$5.97\,\mathrm{s}$) into the full controller ($88\%$/$7.37\,\mathrm{s}$):
$+26$ points of success for $+1.4\,\mathrm{s}$, which shows that the heuristics,
not the base solver, are the source of robustness.

\subsection{Speed/robustness trade-off}
The ablation shows that the geometric rules improve the underlying trade-off
between success rate and lap time, rather than merely exchanging one for the
other. A natural follow-up is whether a single scalar speed parameter, the
cruise target $V_{\mathrm{tgt}}$, could achieve the same gains on its own.
Table~\ref{tab:sweep} sweeps the cruise target $V_{\mathrm{tgt}}$ of the full
controller. Lowering it to $2.0\,\mathrm{m/s}$ trades $0.7\,\mathrm{s}$ of lap
time for $2$ points of success, i.e.\ it stays on the same Pareto front. Pushing
it past the matched value, however, falls off a cliff: at $3.0\,\mathrm{m/s}$
success already drops to $79\%$, and at $3.5\,\mathrm{m/s}$ (well above the soft
per-axis velocity bound $v_{\max}=2.0\,\mathrm{m/s}$) the optimiser flies
permanently against that bound, oversteering into the frames, yet buys almost no
extra time ($6.83\!\to\!6.80\,\mathrm{s}$) while success collapses to $74\%$
(collisions rise to $13\%$). The knob therefore only \emph{slides} the controller
along a success/time Pareto front and quickly saturates; it cannot move the front
itself. The best trade is obtained by matching $V_{\mathrm{tgt}}$ to the
actuation limit, which is why we keep $V_{\mathrm{tgt}}=2.5\,\mathrm{m/s}$. The
identical monotone shape on the held-out Level~3 (Table~\ref{tab:sweep}) shows
this corner-limited behaviour is not an artefact of the single development track.

\begin{table}[t]
\centering
\caption{Cruise-target sweep for the full controller on both levels ($n=100$).}
\label{tab:sweep}
\begin{tabular}{lcccc}
\toprule
& \multicolumn{2}{c}{Level~2} & \multicolumn{2}{c}{Level~3}\\
\cmidrule(lr){2-3}\cmidrule(lr){4-5}
$V_{\mathrm{tgt}}$ [m/s] & Success & Time [s] & Success & Time [s]\\
\midrule
$2.0$ & $90.0\%$ & $8.08$ & $94.0\%$ & $9.17$\\
$2.5$ (default) & $88.0\%$ & $7.37$ & $87.0\%$ & $8.51$\\
$3.0$ & $79.0\%$ & $6.83$ & $83.0\%$ & $7.95$\\
$3.5$ & $74.0\%$ & $6.80$ & $72.0\%$ & $7.74$\\
\bottomrule
\end{tabular}
\end{table}

\subsection{Generalisation to unseen tracks}
Because Level~2 is a single fixed track on which the controller was tuned, we
repeat the ablation on the fully randomised Level~3 as a held-out test
(Fig.~\ref{fig:ablation3} and Table~\ref{tab:ablation3}). The verdict is nuanced. The full adaptation stack
still clearly beats the raw MPCC on unseen tracks ($87\%$ vs.\ $81\%$, collisions
$17\!\to\!10$), so the layer generalises \emph{as a whole}. Per component,
however, only the \emph{curvature cap} transfers cleanly: removing it costs
$4$ points on Level~3 (it was neutral on Level~2), matching the intuition that
random layouts contain sharper corners where the lateral-acceleration limit
actually binds. In contrast, the \emph{gate-track boost} and \emph{caution mode}
(the two components whose radii were tuned to Level~2's fixed reversal and gate
spacing) do not help on Level~3; removing either even nudges success up ($+4$ and
$+3$ points), i.e.\ they are mildly specialised to the development geometry. We
keep them because they are strongly positive on the realistic fixed-track setting
the controller is ultimately graded on, and are not harmful in aggregate, but
this is a clear pointer to make them geometry-adaptive, a natural fit for the
learned weight policy that motivated them.

\begin{table}[t]
\centering
\caption{Held-out ablation on the randomised Level~3 ($n=100$, same seed as
Level~2 so the tracks are paired).}
\label{tab:ablation3}
\begin{tabular}{lccc}
\toprule
Configuration & Success & Mean [s] & Coll.\\
\midrule
Full (all on)          & $87.0\%$ & $8.51$ & $10$\\
$-$ Gate-track boost   & $91.0\%$ & $8.32$ & $5$\\
$-$ Curvature cap      & $83.0\%$ & $7.86$ & $12$\\
$-$ Caution mode       & $90.0\%$ & $7.46$ & $8$\\
All off (raw MPCC)     & $81.0\%$ & $7.30$ & $17$\\
\bottomrule
\end{tabular}
\end{table}
\begin{figure}[!ht]
  \centering
  \includegraphics[width=\columnwidth]{figures/ablation_bars_l3.pdf}
  \caption{Held-out ablation on the randomised Level~3 ($n=100$): success rate
  (bars, left) and mean lap time (line, right). Unlike Level~2
  (Fig.~\ref{fig:ablation}), removing the gate-track boost or caution mode does
  not lower success (these rules are specialised to the fixed track), whereas
  removing the curvature cap does, and the raw controller is still worst.}
  \label{fig:ablation3}
\end{figure}

\subsection{Failure modes}
Table~\ref{tab:crash} breaks the failures down by the nearest physical object at
the moment of the crash, aggregated over the full-controller runs of both levels.
Gate-frame clips dominate on both tracks and concentrate on gate~2 (the tight
$180^\circ$ reversal) and the final gate, where the true pose is revealed latest
and the drone has the least room to correct. Obstacle hits are the second source
and grow slightly under the wider lines the planner draws on Level~3. Ground and
out-of-bounds terminations are rare. This confirms that the limiting factor is
\emph{contouring precision at speed in corners}, not straight-line velocity:
the geometric heuristics that raise precision (boost) or reduce corner entry
speed (curvature cap, caution) each attack this dominant failure mode, which is
why they compound so strongly in the ablation of Table~\ref{tab:ablation}.

\begin{table}[t]
\centering
\caption{Failure breakdown for the full controller ($n=100$ per level).}
\label{tab:crash}
\begin{tabular}{lccccc}
\toprule
Level & Gate frame & Obstacle & Ground & OOB & Timeout\\
\midrule
Level~2 & $4$ & $4$ & $3$ & $0$ & $4$\\
Level~3 & $9$ & $3$ & $0$ & $2$ & $1$\\
\bottomrule
\end{tabular}
\end{table}

\subsection{Discussion}
Across every configuration the dominant failure is a gate-frame clip at the tight
reversal (Table~\ref{tab:crash}), i.e.\ the task is \emph{corner}-limited, not
velocity-limited: the scalar cruise target only slides the controller along a
Pareto front (Table~\ref{tab:sweep}), whereas the geometric rules
(Table~\ref{tab:ablation}) move that front outward. On held-out random tracks the
layer still improves robustness overall, though the two geometry-tuned rules lose
their edge (Table~\ref{tab:ablation3}). This motivates the future directions in
Section~\ref{sec:conclusion}.


\section{Conclusion}
\label{sec:conclusion}
We presented an MPCC-based drone-racing controller that embeds a
collision-aware racing line into the prediction model as a function of a progress
state, eliminating the reference-runaway failure of time-indexed reference
tracking. Three lightweight context-adaptation rules (a curvature speed cap, a
near-gate contouring boost, and a caution mode) lift a raw $62\%$ success rate
to $88\%$ on a fixed track while retaining $87\%$ on fully randomised tracks, and
an ablation study attributes a distinct role to each. A held-out ablation on the
randomised tracks further shows that the physics-based curvature rule generalises,
whereas the two rules whose radii were tuned to the fixed track are mildly
specialised to it, while the stack as a whole still beats the raw controller on
unseen geometry. Deployed unchanged on the physical Crazyflie, the controller
completed all eight test flights at a $8.70\,\mathrm{s}$ mean lap time, within
$0.2\,\mathrm{s}$ of the simulated Level-3 result, demonstrating that the
approach transfers to hardware with only a small sim-to-real gap.

The approach has limitations. Performance is corner-limited: reaching
sub-$5\,\mathrm{s}$ laps would require a structurally better racing line (a
smoother reversal with lower curvature admits a higher cornering speed at the
same lateral-acceleration budget) rather than higher cruise speed. The
adaptation rules, though effective, are hand-designed; the reinforcement-learning
weight scheduler that motivated them could not be trained affordably because the
embedded \texttt{acados} solver is sequential and non-vectorisable, so its online
promise remains untested. Future work therefore targets a minimum-curvature
racing line and completion of the RL weight policy on parallel (Linux/GPU)
infrastructure to make the geometry-specific rules adaptive rather than
hand-tuned.


\bibliographystyle{IEEEtran}
\bibliography{lib}

\end{document}
